\section{System Model}
\label{sec:model}

\begin{figure}[!t]
\centering
\includegraphics[width=\columnwidth]{system.pdf}
\caption{Problem setting. Vibration excitation $v(t)$ drives fretting of
the DLC contacts and modulates their interruption rate. The tool observes
only its own telemetry, from which ECTA also recovers an excitation proxy
$\hat v(t)$.}
\label{fig:system}
\end{figure}

Fig.~\ref{fig:system} summarizes the setting. A diagnostic tool is
attached to the DLC. Through pins 6 and 14 it observes the vehicle's
diagnostic CAN; through pin 16 (battery) and pins 4/5 (ground) it is
powered. We distinguish three tool classes, which differ in what they
record and therefore in what a degraded contact can do to their data:
\begin{itemize}
\item \textbf{Full-bus passive logger}: records every frame on the bus with
  microsecond timestamps (e.g., a SocketCAN interface in listen-only mode).
\item \textbf{Sampling passive logger}: a low-cost device that records only
  the frames it can service (about 65~frames/s in the CANmodes
  hardware, against several hundred to a few thousand on the bus).
\item \textbf{OBD-II polling tester}: requests SAE~J1979 Mode-01 parameters
  over ISO~15765-4~\cite{j1979,iso15765} and records the responses, as
  ELM327-class dongles and fleet telematics units do.
\end{itemize}

\subsection{Wear State}
Let $N(t)$ denote accumulated fretting exposure. Following the
incubation-then-rise behaviour of tin-plated contacts~\cite{antler1985,
bryant1994,park2008}, we model the contact resistance as a logistic law in
log-resistance,
\begin{equation}
\ln R(N)=\ln R_0+\left(\ln R_{\max}-\ln R_0\right)\,
S\!\left(\frac{N-N_{50}}{w}\right),
\label{eq:wear}
\end{equation}
where $S(x)=1/(1+e^{-x})$ is the logistic function, $R_0$ the as-new
resistance, $R_{\max}$ the saturated resistance, $N_{50}$ the
unit-specific mid-life exposure and $w$ the knee width. Exposure
accumulates from vibration and thermal cycling,
\begin{equation}
N(t)=\int_0^t \bar v(\tau)\,d\tau + c_{\mathrm{th}}\, n_{\mathrm{trip}}(t),
\label{eq:exposure}
\end{equation}
with $\bar v$ the deterministic part of the excitation defined below and
$n_{\mathrm{trip}}$ the number of ignition (thermal) cycles. For the
detection experiments we sample $R$ directly from four stages:
healthy ($2$--$50$~m$\Omega$), incipient ($0.3$--$1.5\,\Omega$), moderate
($1.5$--$6\,\Omega$) and severe ($6$--$30\,\Omega$), log-uniformly within
each stage.

\subsection{Excitation}
The vibration intensity at the DLC is modelled as an engine-order term and
a road term, modulated by an unobservable road-input factor:
\begin{equation}
v(t)=\underbrace{\Big(w_e\tfrac{n(t)}{2000}+w_r\tfrac{s(t)}{60}\Big)}_{\bar v(t)}
\exp\!\Big(\sigma_u z(t)-\tfrac{\sigma_u^2}{2}\Big),
\label{eq:vib}
\end{equation}
where $n$ is engine speed (rpm), $s$ is vehicle speed (km/h), and $z(t)$
is a unit Ornstein--Uhlenbeck process with correlation time $\tau_u$.
Crucially, $n(t)$ and $s(t)$ are the \emph{real} signals recorded in each
capture, so every simulated interruption is driven by that vehicle's own
measured operating history, while $z(t)$ ensures that the detector's
excitation proxy is only an imperfect measurement of the true stimulus.

\subsection{Intermittency}
Under vibration, the instantaneous resistance of a fretted contact
fluctuates about $R$, producing the vibration-induced intermittent faults
measured by Shen \emph{et al.}~\cite{shen2018tcpmt}, whose onset shows the
threshold behaviour reported by Flowers \emph{et al.}~\cite{flowers2004}. Modelling the fluctuation as log-normal with spread
$\sigma$, the probability that a micro-motion cycle produces an excursion
above the discontinuity threshold $R_{\mathrm{th}}$~\cite{uscar2} is
$\Phi\big((\ln R-\ln R_{\mathrm{th}})/\sigma\big)$. Interruptions therefore
form an inhomogeneous Poisson process with intensity
\begin{equation}
\lambda(t)=\nu_0\, v(t)\,
\Phi\!\left(\frac{\ln R-\ln R_{\mathrm{th}}}{\sigma}\right),
\label{eq:rate}
\end{equation}
where $\nu_0$ is the rate of micro-motion opportunities at unit
excitation. Interruption durations are log-normal with median
$d_0 (R/R_{\mathrm{th}})^{\gamma}$ and spread $\sigma_d$, capped at
$d_{\max}$; a fraction $p_{\mathrm{pwr}}$ occurs on the supply pins. Since
(\ref{eq:rate}) is proportional to $v(t)$, the loss process inherits the
excitation's temporal structure: this is the signature ECTA exploits.

\subsection{Telemetry Effects}
\emph{Passive tool.} An interruption on CAN-H/L destroys, for the tool
only, every frame whose on-wire interval $[t_i-T_f(i),\,t_i]$ intersects
it, with $T_f=1.1(47+8\,\mathrm{DLC})/\text{bitrate}$. Because the rest of
the network received the frame, nothing is retransmitted: the frame is
simply missing from the log. We assume a listen-only tool, as passive
loggers are; an active tool would also emit error flags that make the fault
visible to the ECUs (e.g., as network trouble codes), an additional
observable that we deliberately do not exploit. A supply interruption longer than the tool's
hold-up time $T_h$ causes a brown-out and silences the tool for
$T_r\sim\mathcal U(0.3,1.5)$~s.

\emph{Polling tester.} We replay each real polling session through the
degraded link, preserving the capture's own inter-transaction gaps. A
transaction whose request or response is on the wire during a CAN-line
interruption is lost and costs the tester a timeout $T_o$; a brown-out
halts polling until re-initialisation completes
($T_r\sim\mathcal U(1,3)$~s), after which the loop resumes. All later
responses shift accordingly, so the simulated session has exactly the
real session's structure except where the fault acts.

\subsection{Confounders}
A useful detector must not confuse DLC wear with other causes of missing
or late data. We model four, each from its own physics:
\begin{enumerate}
\item \textbf{ECU-side intermittent}: the same contact physics at one
  ECU's connector. Its corrupted frames are error-signalled and
  retransmitted after reconnection (late, not lost); interruptions longer
  than \SI{10}{\milli\second} drive the ECU bus-off and silence its
  frames during recovery. Only that ECU's identifiers are affected. In
  polling mode, a supply interruption reboots the engine ECU
  ($0.2$--$0.8$~s), and the tester accumulates consecutive timeouts.
\item \textbf{Bus-wide EMI}: bursts of errors seen by every node,
  including the tool; affected frames are retransmitted back-to-back after
  the burst. Half of the instances are made ignition-coupled (rate
  proportional to engine speed), so that excitation coherence alone cannot
  separate them from wear.
\item \textbf{Logger overflow} (passive): frames beyond the tool's service
  capacity within a \SI{10}{\milli\second} interval are dropped---a
  load-driven, vibration-independent loss.
\item \textbf{Vibration-independent loss} (polling): transactions lost at
  random (ECU busy, tester scheduling), at rates spanning those of
  moderate wear.
\end{enumerate}

\begin{table}[!t]
\caption{Nominal Parameters of the Degradation Model and Their Stress-Test Ranges}
\label{tab:params}
\centering
\footnotesize
\setlength{\tabcolsep}{3pt}
\resizebox{\columnwidth}{!}{%
\begin{tabular}{@{}llll@{}}
\toprule
Symbol & Meaning & Nominal & Stress test (Sec.~\ref{sec:robust})\\
\midrule
$R_{\mathrm{th}}$ & discontinuity threshold & \SI{7}{\ohm}~\cite{uscar2} & 3, \SI{15}{\ohm}\\
$\sigma$ & resistance fluctuation spread & 0.8 & 0.5, 1.2\\
$\nu_0$ & opportunities at $v{=}1$ & \SI{20}{\per\second} & $\times0.5$, $\times2$\\
$d_0,\gamma,\sigma_d$ & interruption duration & \SI{0.2}{\milli\second}, 1, 1 & $d_0\times\frac14$, $\times4$\\
$p_{\mathrm{pwr}}$ & share on supply pins & 0.3 & 0.1, 0.6\\
$T_h$ & tool hold-up time & \SI{2}{\milli\second} & 0.5, \SI{10}{\milli\second}\\
$T_o$ & tester timeout & \SI{0.2}{\second} & ---\\
$w_e, w_r$ & engine / road weights & 0.5, 0.5 & (1,0), (0,1)\\
$\sigma_u,\tau_u$ & unobserved road input & 0.5, \SI{2}{\second} & $\sigma_u{=}0$, 1.0\\
\bottomrule
\end{tabular}}
\end{table}

Table~\ref{tab:params} lists the nominal parameters. Those not fixed by a
standard are unmeasured for J1962 receptacles; rather than claim values for
them, Section~\ref{sec:robust} trains on the nominal model and tests
against each perturbed model.
